\newpage
\phantomsection
\label{prf:closure-under-formation-propositional-formulas}
\LRAProofFor{lem:closure-under-formation-propositional-formulas}

\begin{remark*}[Return]
\hyperref[lem:closure-under-formation-propositional-formulas]{Return to Theorem}
\end{remark*}

\ProofVaultURL{https://github.com/wsollers/lra-proof-vault/tree/master/volume-i/book-logic/propositional-logic/lem-closure-under-formation-propositional-formulas}

\begin{lemma*}[Closure Under Formation]
The set \(\WFF\) is closed under the propositional formation rules.

That is, if \(\varphi\in\WFF\), then
\[
\neg\varphi\in\WFF.
\]
If \(\varphi,\psi\in\WFF\) and
\[
\circ\in\{\land,\lor,\to,\leftrightarrow\},
\]
then
\[
(\varphi\circ\psi)\in\WFF.
\]
\end{lemma*}

\begin{proof}[Professional Standard Proof]
\LRAProofBodyStart
By the recursive definition of \(\WFF\), the negation formation rule says
that whenever \(\varphi\in\WFF\), the formula \(\neg\varphi\) also belongs
to \(\WFF\). Hence
\[
\varphi\in\WFF
\quad\Longrightarrow\quad
\neg\varphi\in\WFF.
\]
The same definition includes the binary connective formation rule: if
\(\varphi,\psi\in\WFF\) and
\[
\circ\in\{\land,\lor,\to,\leftrightarrow\},
\]
then the composite formula \((\varphi\circ\psi)\) belongs to \(\WFF\).
These are exactly the two closure assertions, so \(\WFF\) is closed under
the propositional formation rules.
\end{proof}

\begin{proof}[Detailed Learning Proof]
\LRAProofBodyStart
We prove each formation rule separately.

First suppose that \(\varphi\in\WFF\). One of the clauses in the definition
of a well-formed formula is the negation clause:
\[
\text{if } \theta\in\WFF,\text{ then } \neg\theta\in\WFF.
\]
Applying this clause with \(\theta=\varphi\) gives
\[
\neg\varphi\in\WFF.
\]

Now suppose that \(\varphi,\psi\in\WFF\), and let
\[
\circ\in\{\land,\lor,\to,\leftrightarrow\}.
\]
Another clause in the definition of \(\WFF\) is the binary connective
clause:
\[
\text{if } \theta,\eta\in\WFF
\text{ and }
\bullet\in\{\land,\lor,\to,\leftrightarrow\},
\text{ then } (\theta\bullet\eta)\in\WFF.
\]
Applying this clause with
\[
\theta=\varphi,\qquad \eta=\psi,\qquad \bullet=\circ
\]
gives
\[
(\varphi\circ\psi)\in\WFF.
\]

Since the formulas \(\varphi\), \(\psi\), and the connective \(\circ\) were
arbitrary subject to the stated hypotheses, both required closure properties
hold.
\end{proof}

\begin{remark*}[Proof structure]
This is a definition-unpacking proof. No induction is needed: the lemma only
uses the forward direction of the formation rules that define \(\WFF\).
\end{remark*}

\begin{dependencies}
\begin{itemize}
  \item \hyperref[def:well-formed-formula]{Well-Formed Formula}
\end{itemize}
\end{dependencies}

\clearpage
